%-----------------------------------------------------------------------
\section{Algorithm Deep Dives and Theoretical Complexity}
\label{sec:algo_deep_dives}
%-----------------------------------------------------------------------

To fully contextualize the empirical results, a rigorous theoretical deconstruction of each algorithm's internal computational complexities is required. The transition to post-quantum structures replaces scalar operations (like elliptic curve point multiplications) with substantial array-based polynomial arithmetic.

\subsection{ML-KEM (FIPS 203): Module-LWE Key Encapsulation}
Machine-Lattice Key Encapsulation Mechanism (ML-KEM), formerly known as CRYSTALS-\kyber{}, operates over the polynomial ring $R_q = \mathbb{Z}_q[X]/(X^{256} + 1)$ with a prime modulus $q = 3329$. Let $k$ denote the module rank. ML-KEM generates a matrix of polynomials $\mathbf{A} \in R_q^{k \times k}$. The secret key is a vector of small polynomials $\mathbf{s} \in R_q^k$, and the public key consists of $\mathbf{A}$ and $\mathbf{t} = \mathbf{A}\mathbf{s} + \mathbf{e}$, where $\mathbf{e}$ is a vector of small error polynomials.

The central performance bottleneck in ML-KEM lies in polynomial multiplication. To circumvent basic $O(N^2)$ convolution complexity, ML-KEM mandates the use of the Number Theoretic Transform (NTT), dropping multiplication complexity to $O(N \log N)$. However, on Cortex-M0 (Class 0) and Cortex-M4 (Class 1/2) embedded devices, executing 256-point NTT algorithms demands substantial register availability. Since $q=3329$ admits a 256-th root of unity $\zeta$, polynomials are decomposed into 128 linear polynomials modulo $X^2 - \zeta^{2i+1}$. While highly optimized in assembly (utilizing M4's DSP extensions like \texttt{SMLAD}), the initial generation of the uniform matrix $\mathbf{A}$ via XOFs (SHAKE-128) consumes up to 40\% of the key generation clock cycles, heavily taxing battery life during the constant symmetric hashing phases.

\subsection{ML-DSA (FIPS 204): Module-LWE Signatures}
ML-DSA (\dilithium{}) operates symmetrically to ML-KEM but adopts the "Fiat-Shamir with Aborts" framework to construct signatures from interactive zero-knowledge proofs. Unlike KEM operations, digital signatures must enforce strict constant-time bounds to thwart timing-based side channels. 

During signing, ML-DSA samples a random masking vector $\mathbf{y}$, calculates an intermediate commitment $\mathbf{w}_1 = \mathbf{A}\mathbf{y}$, and hashes it with the message to produce a challenge $c$. The prover computes the potential signature $\mathbf{z} = \mathbf{y} + c\mathbf{s}_1$. Crucially, to prevent secret key leakage through $c\mathbf{s}_1$, the signature is \textit{rejected} and restarted if the coefficients of $\mathbf{z}$ exceed strict predefined bounds (rejection sampling).

In resource-constrained environments, this rejection sampling is devastating. The non-deterministic loop nature implies that average-case clock cycles heavily diverge from best-case paths. For Class 1 devices lacking large CPU caches, continuously regenerating large polynomial vectors and re-evaluating NTTs upon every rejection abort drains energy reserves exponentially faster than classical ECDSA deterministic loops. Furthermore, ML-DSA public keys and signatures scale upward of 2.5 KB (ML-DSA-44), demanding fragmented DTLS transmission buffers.

\subsection{SLH-DSA (FIPS 205): Stateless Hash-Based Signatures}
SLH-DSA (SPHINCS+) represents the most conservative cryptographic safety profile, relying exclusively on the quantum resistance of the underlying hash functions (e.g., SHA-256 or SHAKE-256). It entirely avoids unproven algebraic constructs like Module-LWE.

The architecture comprises a substantial hyper-tree (a tree of Merkle trees). A signature authenticates a message by signing it with a Few-Time Signature (FORS) scheme at the bottom leaf layer, and then providing the full authentication path traversing up exactly $d$ layers of Merkle trees to the root public key. 

While mathematically pristine, SLH-DSA's performance profile on IoT endpoints is uniformly hostile. Constructing a single signature requires executing tens of thousands of hash compression functions. More critically, the resulting signature encapsulates the entire authentication path. An SLH-DSA-128s signature consumes roughly 7,856 bytes. Transmitting an 8 KB signature over a standard IEEE 802.15.4 LPWAN link (with a 127-byte Maximum Transmission Unit) requires over 80 discrete radio frames. If a single frame drops, the entire signature fails verification, rendering SLH-DSA virtually impossible to deploy robustly on Class 1 UDP/CoAP constrained endpoints without extensive error correction protocols.

\subsection{FN-DSA (Draft FIPS 206): \falcon{} and \ntru{} Lattices}
Fast-Fourier Lattice-Based Compact Signatures over \ntru{} (\falcon{}) solves the substantial signature size dilemma of ML-DSA (reducing sizes down to 666 bytes for 128-bit security). It achieves this by utilizing extremely dense \ntru{} lattices where the public key is merely a single polynomial $h = g f^{-1} \bmod q$.

However, to generate signatures, FN-DSA uses "Hash-and-Sign" over \ntru{} lattices guided by a Fast Fourier Sampling algorithm across a continuous Gaussian distribution. Executing true Gaussian sampling securely is famously complex. Specifically, \falcon{} requires high-precision 64-bit IEEE-754 floating-point arithmetic. 

For IoT environments, especially Class 0 and Class 1 ARM Cortex processors frequently lacking a hardware Floating Point Unit (FPU), these operations must be emulated entirely in software. Software FPU emulation invokes thousands of instruction cycles for basic multiplication and division boundaries, resulting in FN-DSA key generation times that routinely exceed 20,000 milliseconds (20 seconds) on 48MHz processors, triggering watchdog timer resets blocking physical deployment.

\subsection{Hamming Quasi-Cyclic (HQC)}
Introduced broadly as an alternative KEM standard to hedge against theoretical breakthroughs in lattice cryptography, HQC relies on coding theory—specifically, the hardness of decoding random quasi-cyclic codes coupled with resolving the syndrome decoding problem.

HQC provides an intriguing trade-off matrix: while the public key sizes are excessively large (2,249 bytes for 128-bit security, compared to 800 bytes for ML-KEM), the actual mathematical operations bypass NTT convolutions entirely, opting for simpler vector additions and multiplications over finite fields mapped into GF(2). For ultra-low power bare-metal Class 0 devices utilizing basic bit-shift architectures, HQC decoding algorithms frequently consume strictly less instantaneous wattage than ML-KEM, despite the heavy radio penalty of receiving substantial ciphertexts.
